\subsection{Poisson Observation Model}

We consider observations that follow Poisson distributions, which is natural for count data such as disease cases:
\begin{equation}
Y_i(t) \sim \text{Poisson}(\lambda_i(t))
\end{equation}
where $\lambda_i(t) = N \cdot \mu_i(t)$ is the expected number of cases in region $i$ at time $t$, with $N$ representing the population scale and $\mu_i(t) = \beta(t) S_i(t) \sum_j C_{ij} I_j(t)$ the mean incidence rate from our SIR model.

The observation vector is $\mathbf{Y} = [Y_1(0), Y_2(0), \ldots, Y_1(T-1), Y_2(T-1)]^T \in \mathbb{N}_0^{2T}$, with mean vector:
\begin{equation}
\boldsymbol{\lambda}(\boldsymbol{\phi}) = [\lambda_1(0), \lambda_2(0), \ldots, \lambda_1(T-1), \lambda_2(T-1)]^T
\end{equation}

\subsubsection{Statistical Properties}

Under the Poisson model:
\begin{align}
\mathbb{E}[Y_i(t)] &= \lambda_i(t) = N \mu_i(t) \\
\text{Var}[Y_i(t)] &= \lambda_i(t) = N \mu_i(t)
\end{align}

The variance equals the mean, which is a characteristic property of Poisson distributions.

\subsubsection{Likelihood Function}

The likelihood function for independent Poisson observations is:
\begin{equation}
f(\mathbf{Y}|\boldsymbol{\phi}) = \prod_{t=0}^{T-1} \prod_{i=1}^{2} \frac{\lambda_i(t)^{Y_i(t)} e^{-\lambda_i(t)}}{Y_i(t)!}
\end{equation}

\subsubsection{Log-Likelihood Function}

Taking the logarithm:
\begin{align}
\ell(\boldsymbol{\phi}) &= \log f(\mathbf{Y}|\boldsymbol{\phi}) \\
&= \sum_{t=0}^{T-1} \sum_{i=1}^{2} [Y_i(t) \log \lambda_i(t) - \lambda_i(t) - \log(Y_i(t)!)] \\
&= \sum_{t=0}^{T-1} \sum_{i=1}^{2} [Y_i(t) \log \lambda_i(t) - \lambda_i(t)] + \text{const}
\end{align}

where the constant term does not depend on the parameters $\boldsymbol{\phi}$.

\subsubsection{Score Functions}

Taking the partial derivative with respect to parameter $\phi_k \in \{\theta, S_1(0), I_1(0), S_2(0), I_2(0)\}$:
\begin{align}
\frac{\partial \ell}{\partial \phi_k} &= \sum_{t=0}^{T-1} \sum_{i=1}^{2} \left[Y_i(t) \frac{1}{\lambda_i(t)} - 1\right] \frac{\partial \lambda_i(t)}{\partial \phi_k} \\
&= \sum_{t=0}^{T-1} \sum_{i=1}^{2} \left[\frac{Y_i(t) - \lambda_i(t)}{\lambda_i(t)}\right] \frac{\partial \lambda_i(t)}{\partial \phi_k}
\end{align}

\subsubsection{Fisher Information Matrix}

Using the outer product of gradients method:
\begin{align}
[\mathbf{J}]_{kl} &= \mathbb{E}\left[\frac{\partial \ell}{\partial \phi_k}\frac{\partial \ell}{\partial \phi_l}\right] \\
&= \mathbb{E}\left[\sum_{t,i} \sum_{s,j} \frac{Y_i(t) - \lambda_i(t)}{\lambda_i(t)} \cdot \frac{Y_j(s) - \lambda_j(s)}{\lambda_j(s)} \cdot \frac{\partial \lambda_i(t)}{\partial \phi_k}\frac{\partial \lambda_j(s)}{\partial \phi_l}\right] \\
&= \sum_{t,i} \sum_{s,j} \mathbb{E}\left[\frac{Y_i(t) - \lambda_i(t)}{\lambda_i(t)} \cdot \frac{Y_j(s) - \lambda_j(s)}{\lambda_j(s)}\right] \frac{\partial \lambda_i(t)}{\partial \phi_k}\frac{\partial \lambda_j(s)}{\partial \phi_l}
\end{align}

Since the Poisson observations are independent across time and regions:
\begin{equation}
\mathbb{E}\left[\frac{Y_i(t) - \lambda_i(t)}{\lambda_i(t)} \cdot \frac{Y_j(s) - \lambda_j(s)}{\lambda_j(s)}\right] = \begin{cases}
\frac{1}{\lambda_i(t)} & \text{if } (i,t) = (j,s) \\
0 & \text{if } (i,t) \neq (j,s)
\end{cases}
\end{equation}

This follows because for a Poisson random variable $Y \sim \text{Poisson}(\lambda)$:
\begin{align}
\mathbb{E}\left[\frac{Y - \lambda}{\lambda}\right] &= 0 \\
\text{Var}\left[\frac{Y - \lambda}{\lambda}\right] &= \frac{\text{Var}[Y]}{\lambda^2} = \frac{\lambda}{\lambda^2} = \frac{1}{\lambda}
\end{align}

Therefore:
\begin{align}
[\mathbf{J}]_{kl} &= \sum_{t=0}^{T-1}\sum_{i=1}^{2} \frac{1}{\lambda_i(t)} \frac{\partial \lambda_i(t)}{\partial \phi_k}\frac{\partial \lambda_i(t)}{\partial \phi_l} \\
&= \sum_{t=0}^{T-1}\sum_{i=1}^{2} \frac{1}{N \mu_i(t)} \frac{\partial (N \mu_i(t))}{\partial \phi_k}\frac{\partial (N \mu_i(t))}{\partial \phi_l} \\
&= \sum_{t=0}^{T-1}\sum_{i=1}^{2} \frac{N}{\mu_i(t)} \frac{\partial \mu_i(t)}{\partial \phi_k}\frac{\partial \mu_i(t)}{\partial \phi_l}
\end{align}

\subsubsection{Matrix Form}

Define the Jacobian matrix:
\begin{equation}
\mathbf{G} = \frac{\partial \boldsymbol{\mu}}{\partial \boldsymbol{\phi}} =  \begin{bmatrix}
\frac{\partial \mu_1(0)}{\partial \theta} & \frac{\partial \mu_1(0)}{\partial S_1(0)} & \cdots & \frac{\partial \mu_1(0)}{\partial I_2(0)} \\
\frac{\partial \mu_2(0)}{\partial \theta} & \frac{\partial \mu_2(0)}{\partial S_1(0)} & \cdots & \frac{\partial \mu_2(0)}{\partial I_2(0)} \\
\vdots & \vdots & \ddots & \vdots \\
\frac{\partial \mu_2(T-1)}{\partial \theta} & \frac{\partial \mu_2(T-1)}{\partial S_1(0)} & \cdots & \frac{\partial \mu_2(T-1)}{\partial I_2(0)}
\end{bmatrix} \in \mathbb{R}^{2T \times 5}
\end{equation}

And the weight matrix:
\begin{equation}
\mathbf{W} = \text{diag}\left(\frac{N}{\mu_1(0)}, \frac{N}{\mu_2(0)}, \ldots, \frac{N}{\mu_2(T-1)}\right) \in \mathbb{R}^{2T \times 2T}
\end{equation}

Then the Fisher Information Matrix can be written compactly as:
\begin{equation}
\mathbf{J} = \mathbf{G}^T \mathbf{W} \mathbf{G} = \sum_{t=0}^{T-1}\sum_{i=1}^{2} \frac{N}{\mu_i(t)} \frac{\partial \mu_i(t)}{\partial \boldsymbol{\phi}} \frac{\partial \mu_i(t)}{\partial \boldsymbol{\phi}}^T
\end{equation}

\subsubsection{Scaling Properties}

The Poisson model exhibits natural scaling with population size $N$:

\begin{itemize}
\item \textbf{Fisher Information}: $\mathbf{J} \propto N$ (increases linearly with population)
\item \textbf{CRLB}: $\text{Var}(\hat{\theta}) \geq [\mathbf{J}^{-1}]_{11} \propto 1/N$ (decreases with population)
\item \textbf{Standard Deviation}: $\text{SD}(\hat{\theta}) \propto 1/\sqrt{N}$ (classical $\sqrt{N}$ convergence)
\end{itemize}

This scaling behavior is much more reasonable than the multiplicative Gaussian model, as larger populations provide more information about the connectivity parameters, but the improvement follows the standard statistical rate of $\sqrt{N}$.

\subsubsection{Cramér-Rao Lower Bound}

The Cramér-Rao lower bound for the connectivity parameter is:
\begin{equation}
\text{Var}(\hat{\theta}) \geq [\mathbf{J}^{-1}]_{11}
\end{equation}

This can be computed numerically using Cholesky factorization as described in previous sections. The Poisson model provides a more realistic assessment of parameter uncertainty compared to pure multiplicative Gaussian noise, especially for small incidence rates where the Gaussian approximation may break down.
